%% ========================================================================
%%%% Layout-Einstellungen (unverändert aus Bericht_Vorlage übernommen)
%% ========================================================================

%% How to separate paragraphs: indention ("no") or spacing ("half", "full", ...).
\newcommand{\myparskip}{half}

%% Inner binding correction.
\newcommand{\myBCOR}{0mm}

%% oneside oder twoside
\newcommand{\mylaterality}{oneside}

%% Farbe der Überschriften etc. in RGB
\newcommand{\mydispositioncolor}{22,72,154}

%% Zeilenabstand
\newcommand{\mylinespread}{1.5}

\newcommand{\mycolorlinks}{true}

%% babel-Sprache (die *letzte* Sprache ist die aktive)
\newcommand{\mylanguage}{american,ngerman}

%% ========================================================================
%%%% Dokument-Metadaten -- Thema 2: KUKA Ready2educate
%% ========================================================================

%% \myauthor wird nur für die PDF-Metadaten (pdfauthor/pdfproducer in
%% pdf_settings.tex) verwendet, nicht auf dem Deckblatt angezeigt.
\newcommand{\myauthor}{Projektteam Thema 2 -- KUKA Ready2educate}

%% Namen und Matrikelnummern
%% Erscheinen auf dem Deckblatt (siehe coverpage.tex).
\newcommand{\myauthors}{%
    Nico Schenkel\\
    Nicolas Mahr\\
    Maike Lindörfer\\
    Maria Fernanda Ledezma%
}
\newcommand{\mymatrikelnummern}{%
    1234567\\
    1407877\\
    1234567\\
    1234567%
}

\newcommand{\myworktype}{Laborbericht}
\newcommand{\myworktitle}{}
\newcommand{\mytitle}{KUKA Ready2educate: Echtzeit-Anbindung und kraftgeregelte Systemintegration}
\newcommand{\mysubtitle}{ }
\newcommand{\mysubject}{Echtzeitfähige Robotersteuerung und Kraftregelung}
\newcommand{\mykeywords}{KUKA, QUARC, Echtzeit, KVP, EtherCAT, Kraftregelung}

\newcommand{\mydegree}{Master of Science}

%% Duale Hochschule / Unternehmen: entfällt, da kein duales Praxisprojekt
%% (siehe Vorlage: \mydualpartner, \mydualpartnerLocation, \mysupervisor
%% werden hier nicht definiert, da coverpage.tex die entsprechenden Zeilen
%% ohnehin auskommentiert hat).

\newcommand{\myuniversity}{Frankfurt University of Applied Sciences}
\newcommand{\mycourse}{Simulation und Regelung (SoSe 2026)}
\newcommand{\mystudy}{Mechatronik und Robotik}
\newcommand{\myevaluator}{K. Schmidt}

%% TODO: Abgabedatum noch offen -- eintragen, sobald bekannt.
\newcommand{\mysubmissionday}{TT}
\newcommand{\mysubmissionmonth}{MM}
\newcommand{\mysubmissionyear}{2026}

\newcommand{\myprocessingperiodbegin}{xxx}
\newcommand{\myprocessingperiodend}{xxx}
